% ---------------------------------------------------------------------------
% The manuscript body, shared verbatim by the Springer (LNCS) and IEEE
% (IEEEtran) builds. Nothing class-specific may appear in this file.
%
% Each wrapper defines, before \input-ing this:
%   \wfigure / \endwfigure   a figure spanning the full text width
%   \wtable  / \endwtable    the same for tables
%   \tabscale                the font size command tables are set at
%   \para{...}               a run-in bold paragraph heading
%   \citep{...}               mapped to \cite for the LNCS/IEEEtran builds
% ---------------------------------------------------------------------------

\section{Introduction}

Building a marine species-distribution or harmful-algal-bloom model, or any
analytics feature that needs occurrence, catch, or effort data broader than
what one platform already fetches, starts with a question that sounds like it
should have a settled answer: which data sources are actually reachable, and
in what shape? In practice, the answer is scattered across API documentation
of wildly uneven quality, mailing-list folklore, a portal's own about page,
and direct trial and error, and it decays --- an endpoint that returned JSON
last quarter can start returning an HTML redirect, a certificate error, or
nothing at all, for reasons that have nothing to do with the data it holds.
Sambasivan et al.\ document, at the level of organisational practice, that
this kind of unglamorous data work is systematically undervalued relative to
model work \citep{sambasivan2021}; a source-reachability survey is exactly
the least glamorous end of that work, which is also why it is rarely written
down, and almost never written down with a date attached.

This paper reports one such survey: \NSourcesSurveyed{} marine
biodiversity, fisheries, and catch/effort data sources, each probed live on
\ProbeDate{} against its documented or discoverable API endpoint. Each is
classified into one of three shapes: open REST/JSON and occurrence-shaped;
reachable but gated or bulk-download-only; or dead and defunct. Each verdict
is reported with the specific HTTP status, DNS or TLS failure, or record
count the probe returned. The sources span the field's usual anchors (OBIS, GBIF), regional
and taxonomic infrastructure (WoRMS, EMODnet Biology, ICES DATRAS), regional
fisheries management organisations (IOTC, ICCAT, WCPFC), and a set of Indian
marine research institutions whose accessibility motivated the survey in the
first place. The paper's own thesis is reflexive: a survey of this kind is
useful precisely because it goes stale, which is the argument for publishing
it now, with dated probes, rather than waiting for a version that will not
arrive. Section~\ref{sec:results} demonstrates this directly rather than only
asserting it: re-probing a subset of the \NSourcesSurveyed{} sources while
writing this paper, more than a month after the original pass, already finds
\NEndpointsDrifted{} of \NEndpointsReprobed{} re-checked endpoints answering
with a different technical signature than they did on \ProbeDate{}.

\subsection{Contributions}

\begin{enumerate}\itemsep2pt
\item A \textbf{dated, three-tier survey} of \NSourcesSurveyed{} marine
biodiversity/fisheries/catch data sources, each with a live-probed verdict,
the specific technical detail behind it (status code, DNS/TLS failure mode,
or record count), and the date it was true.
\item Three \textbf{corrections to a plausible-sounding default}, each
demonstrated with the specific measurement rather than asserted: GBIF is not
a superset of OBIS for marine fish, so unioning the two requires
deduplication rather than a source swap; a single non-200 probe response
does not by itself mean a source is dead, since the same class of response
(status~000) meant a transient TLS handshake flake for one source and a
genuinely unreachable host for another; and an institution's parent-domain
hostname failing does not mean the institution is unreachable, since two of
four Indian institutes sharing a dead \texttt{icar.gov.in} subdomain are
reachable under their own \texttt{.res.in} domain instead.
\item An explicit \textbf{evidence-grading stance}: every claim in this paper
is graded \emph{Dated} rather than \emph{Artifact}, because the object under
test is infrastructure this project does not control, and Section
\ref{sec:method} explains why no claim here could honestly be graded
otherwise.
\item A \textbf{live re-verification pass}, performed while writing this
paper, that finds the survey's own numbers already drifting in a way that
argues for the paper's central claim rather than against it: an endpoint
recorded dead by DNS failure on \ProbeDate{} now resolves to a live host
serving an unconfigured default certificate, and a portal recorded as 404
now redirects to a live page while its programmatic API remains unchanged.
\item The survey's \textbf{sharpest single finding}, isolated in
Section~\ref{sec:discussion}: no Indian marine data source runs its own
machine-readable API --- except that this is not quite true of IndOBIS, which
needs no new integration at all, because it is already reachable as a node
filter on the OBIS API this codebase already calls.
\end{enumerate}

\section{Related Work}
\label{sec:related}

\para{What the major aggregators' own papers describe} Grassle's founding
description of the Ocean Biogeographic Information System \citep{grassle2000}
and Robertson et al.'s account of the GBIF Integrated Publishing Toolkit
\citep{robertson2014} both describe what these two aggregators are, how data
enters them, and why they were built as shared infrastructure rather than
per-project scraping. Costello et al.\ give the equivalent account for the
World Register of Marine Species, the taxonomic backbone this survey lists
as the right place to resolve a name before querying anything else
\citep{costello2013worms}. None of the three asks, or is designed to answer,
whether a given consumer can reach a given endpoint programmatically today,
nor how a source outside this small, well-resourced set behaves --- the
question this survey's \NTierTwo{}-source gated/bulk tier and
\NTierThree{}-source Indian-institution tier exist to answer, since neither
tier is the subject either infrastructure paper describes.

\para{The closest prior assessment, and where this survey sits upstream of
it} Vandepitte et al.\ assess the data quality, completeness, and fitness for
use of records already inside marine biogeographic databases --- duplicate
records, taxonomic mismatches, geographic outliers --- once a dataset has
already been published into an aggregator like OBIS \citep{vandepitte2015}.
That is a rigorous, valuable question, and it is a different one from this
survey's: Vandepitte et al.\ examine data already reachable and already
inside the system this survey's Tier 1 describes, whereas this survey asks
the question upstream of that --- is a given source reachable and shaped as
advertised at all --- across a much longer tail of sources, most of which
have never been the subject of a quality assessment because nobody has
systematically probed whether they answer a query in the first place.
Neither paper substitutes for the other; a source has to clear this survey's
bar before Vandepitte et al.'s questions can even be asked of it.

\para{Scope} This is a probe-based accessibility survey, not a data-quality
audit and not a taxonomic or geographic coverage assessment of the sources it
lists; \S\ref{sec:limitations} states directly that its own findings are
exactly the kind of claim this paper argues decays fastest.

\section{Method: what \emph{probed} means, and why every claim here is
\emph{Dated}}
\label{sec:method}

\para{Probing} Each source in Table~\ref{tab:open}, Table~\ref{tab:gated},
and Table~\ref{tab:indian} was queried live against its documented or
discoverable API root or a representative endpoint (an occurrence search, a
statistics call, a WFS \texttt{GetCapabilities} request, or the plainest
public page where no API is documented at all), recording the returned HTTP
status code, DNS or TLS failure where the request did not complete, and,
where the response was 200, whether the body was machine-readable (JSON,
XML, or WFS) and occurrence-shaped, a generic repository schema, or plain
HTML. No source was probed through an authenticated session beyond a key
already provisioned in this codebase's own configuration
(\texttt{backend/.env}); a 401/403 response is reported as \emph{gated}, not
retried with credentials this project does not hold. The original pass was
run \ProbeDate{} and recorded in \texttt{docs/data-sources.md} (split out of
\texttt{CLAUDE.md} on \SplitDate{} to keep the root instruction file
lean); this paper adds a second, partial pass of \NEndpointsReprobed{}
endpoints, re-run \AsOf{} while writing it, reported in
\S\ref{sec:results}.

\para{The three-tier scheme} \emph{Tier 1} (\NTierOne{} sources, Table
\ref{tab:open}): open REST/JSON or equivalent (WFS, OAI-PMH) with no key
required for reads, returning occurrence- or dataset-shaped data. \emph{Tier
2} (\NTierTwo{} sources): reachable in some form --- a documentation page, a
portal, a bulk-download form --- but gated behind a key or account, or usable
only as a file download rather than a query. \emph{Tier 3} (\NTierThree{}
sources): Indian marine research institutions, considered separately because
they motivated the survey and because, as \S\ref{sec:discussion} argues, they
share a specific and previously undocumented pattern that a plain
Tier-1/Tier-2 split would flatten.

\para{Why every claim here is graded Dated, not Artifact} A companion
catalogue from the same research programme grades its own claims
\emph{Artifact} when a committed report file or a currently-passing
regression test reproduces the number on demand, and \emph{Dated} when the
only reproduction path is a named script run again against a moving
field --- a day's wind, current, or SST state
\citep{silentfailures2026}\footnote{Cross-referenced for its grading
convention only; this survey has no numerical or mechanistic overlap with
that catalogue's eleven instances.}. Every claim in this paper falls further
along that same axis than either of that catalogue's own categories: there is
no script this project owns whose re-run reproduces a source's HTTP status,
because the object being measured is infrastructure entirely outside this
project's control. Re-running the same request tomorrow is not
\emph{re-computation}, the way re-running a training script against a
committed feature store is; it is a fresh observation of a system that can
legitimately answer differently for reasons --- a certificate renewal, a
migrated domain, a maintenance window, a policy change --- that have nothing
to do with anything this project did. Accordingly, this paper does not use an
\emph{Artifact} grade at all: every verdict below is \emph{Dated}, stamped
with the date it was observed, and \S\ref{sec:results}'s re-verification
pass is offered as a demonstration of exactly how much that distinction
matters over as little as five weeks.

\section{Results}
\label{sec:results}

\begin{wtable}
\centering
\tabscale
\caption{Tier 1: open REST/JSON (or equivalent) sources, no key required for
reads, occurrence- or dataset-shaped. All probed \ProbeDate{} unless a
re-probe date is noted; \NTableRowsOpen{} rows cover \NTierOne{} sources.}
\label{tab:open}
\begin{tabular}{p{0.14\linewidth}p{0.08\linewidth}p{0.70\linewidth}}
\toprule
Source & Status & Detail \\
\midrule
OBIS & 200 & \texttt{/v3/occurrence}, \texttt{/statistics}, \texttt{/checklist} all 200; \texttt{hasextensions=DNADerivedData} eDNA filter returns \ObisEdnaRecordsM{}M records, \ObisEdnaSpecies{} species, \ObisEdnaDatasets{} datasets, 2001--2025. Already implemented (\texttt{marine\_ml/sources/obis.py}). \\
GBIF & 200 & \texttt{/v1/occurrence/search}, no key, Darwin Core. Already implemented (\texttt{sources/gbif.py}); see below for why unioning it with OBIS is not a superset relation. \\
iNaturalist & 200 & \texttt{/v1/observations}, no key; citizen-science coastal bias, \texttt{quality\_grade=research} mandatory; not a survey. \\
WoRMS & 200 & \texttt{/rest/AphiaRecordsByName}; taxonomic backbone for resolving a name/AphiaID before querying anything else. \\
Marine Regions & 200 & \texttt{/rest/getGazetteerRecordsByName.json}; EEZ/LME/sea-area polygons for defining a region bound instead of a hand-written bbox. \\
Sea Around Us & 200 & \texttt{/api/v1/eez}, \texttt{/eez/tonnage/taxon}; reconstructed catch by EEZ/taxon --- a different table from occurrences. \\
ICES DATRAS & 200 & \texttt{/getSurveyList} (XML, SOAP-style \texttt{.asmx}); the only source in this table with true zero-catch hauls, i.e.\ true absences; North Atlantic/European coverage only. \\
EMODnet Biology & 200 & WFS \texttt{GetCapabilities}; \texttt{GetFeature}+bbox works without a portal scrape; European coverage. \\
NASA Earthdata CMR & 200 & \texttt{/search/collections.json} unauthenticated for search; downloads require Earthdata Login. \\
NOAA ERDDAP & 200 & Known in this codebase to flap --- unloads and reloads datasets, answering a transient 404 indistinguishable from a removed dataset (applies to every ERDDAP host, OTN included). \\
OTN & 200 & Acoustic telemetry (individual tracks, not occurrences); its \texttt{/info/index.json} probe alone returned 000, a TLS handshake flake, not treated as a dead source --- see the correction below. \\
HYCOM & 200 & THREDDS/OPeNDAP; environmental covariate, overlaps Copernicus, no reason to add unless a field is missing there. \\
Zenodo / Dryad / DataONE & 200 & Generic repositories; per-dataset ad-hoc schema, useful for discovering a dataset, not for a scheduled fetch. \\
PANGAEA & 200 / 404 & OAI-PMH provider (\texttt{?verb=Identify}) 200, re-probed \AsOf{}: unchanged. Its older \texttt{/es/dataportal-item/\_search} endpoint is 404 --- gone, use OAI-PMH instead --- re-probed \AsOf{}: also unchanged. \\
\bottomrule
\end{tabular}
\end{wtable}

\begin{wtable}
\centering
\tabscale
\caption{Tier 2: reachable but gated or bulk-download only. All probed
\ProbeDate{} unless a re-probe date is noted; \NTableRowsGated{} rows cover
\NTierTwo{} sources.}
\label{tab:gated}
\begin{tabular}{p{0.16\linewidth}p{0.10\linewidth}p{0.66\linewidth}}
\toprule
Source & Status & Detail \\
\midrule
Global Fishing Watch & 401 & Expected; needs the token already provisioned in \texttt{backend/.env} (\texttt{services/gfw.py}). Effort/AIS, not biology. \\
Movebank & 401 & Free account plus per-study permission from the data owner; most studies are not open. \\
AIMS & 200 / 403 & Docs 200; \texttt{api.aims.gov.au} 403 --- needs an API key. \\
FishBase & 000 & \texttt{fishbase.ropensci.org} dead by DNS on \ProbeDate{}; re-probed \AsOf{}, its failure signature has changed --- see below. \texttt{fishbase.se} serves 200 HTML only; programmatic use means the \texttt{rfishbase} data dumps, not a live API. \\
AquaMaps / Bio-ORACLE / MARSPEC & 200 & Modelled raster products downloaded as files, not queried; AquaMaps output is itself a model prediction, so using it as a label trains on another model. \\
IOTC / ICCAT / CCSBT & 200 & RFMO catch/effort, gridded aggregates behind manual download forms; ICCAT partly requires a request. \\
WCPFC & 404 & Both \texttt{/public-domain} and \texttt{/public-domain-data} 404 on \ProbeDate{} --- the URL moved; re-probed \AsOf{}: still 404. \\
FAO FishStat & 200 & Desktop application plus bulk zip; country-year granularity, fine for context, not for a habitat model. \\
NOAA Stock SMART & 200 / 404 & Web app 200; \texttt{/sis/api/} and \texttt{/adb/} both 404 --- no documented public API, treat as bulk export. \\
DFO Canada & 200 / 400 & Federal CKAN \texttt{package\_search} 200; the plain \texttt{/dataset?organization=} URL 400s. \\
CSIRO / NCEI / MBON / GOOS & 200 & Reachable only as portals. \\
NIWA & 200 / 404 & Root 200; \texttt{/coasts-and-oceans/marine-data} 404. \\
RLS & 200 & Standardised reef transects with abundance and real zeros --- genuinely valuable --- but ships as CSV downloads only. \\
VertNet & 404 / 500 & Portal 404, appspot API 500 on \ProbeDate{}; re-probed \AsOf{}, the portal's status has changed but the API has not --- see below. \\
\bottomrule
\end{tabular}
\end{wtable}

\begin{wtable}
\centering
\tabscale
\caption{Tier 3: Indian marine research institutions, no machine-readable
API (one important exception). All probed \ProbeDate{} unless a re-probe
date is noted; \NTableRowsIndian{} rows cover \NTierThree{} sources.}
\label{tab:indian}
\begin{tabular}{p{0.16\linewidth}p{0.10\linewidth}p{0.66\linewidth}}
\toprule
Source & Status & Detail \\
\midrule
CMFRI & 200 & \texttt{cmfri.org.in} and \texttt{eprints.cmfri.org.in} both 200; the EPrints repository has OAI-PMH and holds annual landings estimates, but document-level, not record-level. \\
FSI & 200 & Survey cruise data; publications only, data by request. \\
IndOBIS & 200 & Not a counterexample to ``no machine-readable API'' in the way it first appears --- reachable through the existing OBIS API as node \texttt{\IndObisNodeId{}}, no new integration needed. See \S\ref{sec:discussion}. \\
CIFT / CIBA & 200 & \texttt{cift.res.in}, \texttt{ciba.res.in} both 200, re-probed \AsOf{}: unchanged; their \texttt{*.icar.gov.in} subdomains are 000. \\
CIFRI / CIFA & 000 / 400 & \texttt{*.icar.gov.in} subdomains 000 (DNS/TLS); \texttt{cifa.nic.in} 400. No confirmed working \texttt{.res.in} equivalent for either, unlike CIFT/CIBA. \\
\bottomrule
\end{tabular}
\end{wtable}

Of \NSourcesSurveyed{} sources probed \ProbeDate{}, \NTierOne{} are open,
occurrence-shaped, and require no key; \NTierTwo{} are reachable in some form
but gated or bulk-download-only; and the \NTierThree{} Indian institutions in
Table~\ref{tab:indian} have, with one important exception discussed below,
no machine-readable API at all --- everything is a PDF or HTML publication,
a portal, or a manual request.

\para{Correction 1: GBIF is not a superset of OBIS for marine fish}
Unioning a second occurrence source with an existing one is often treated as
a strict superset operation --- more of the same kind of data, so nothing is
lost by adding it. Measured directly on \emph{Thunnus albacares}, OBIS
returns \ObisTunaRecords{} records against GBIF's \GbifTunaRecords{} with
coordinates: a large fraction of GBIF's marine records are themselves OBIS
datasets republished, so a naive union double-counts, and the two sources
dominate for different species groups in this codebase's own target list ---
GBIF for the pelagic tunas, OBIS by roughly an order of magnitude for Indian
mackerel and oil sardine, the two species that matter most for Indian
coastal fisheries. The union is a genuine union, not a migration, and it
requires deduplication on \texttt{occurrenceID} (falling back to dataset,
catalogue number, latitude/longitude, and date) precisely because neither
source can be dropped for the other.

\para{Correction 2: a single non-200 probe is not, by itself, a death
certificate} Two sources across Table~\ref{tab:open} and
Table~\ref{tab:gated} returned status 000 --- a connection that did not
complete at all --- on a single probe, and only one of them is actually dead.
OTN's \texttt{/info/index.json} endpoint returned 000 on \ProbeDate{}, and a
second request to the same host's plain HTML index succeeded; read alone,
the 000 would have supported the same ``dead'' verdict this survey gives
FishBase's ropensci mirror, whose 000 was a genuine DNS resolution failure
that did not resolve on a retry. The methodological point this survey's own
working notes make explicit is that a status-000 probe needs a second,
different request against the same host before it is read as evidence of
death rather than of a transient handshake flake --- exactly the same
discipline a companion catalogue from this codebase applies to a corroboration
rate needing a control arm before it is read as confirmation
\citep{silentfailures2026}.

\para{Correction 3: a dead parent domain does not mean a dead institution}
Four Indian institutes --- CIFRI, CIFT, CIBA, and CIFA --- share a hostname
pattern under \texttt{icar.gov.in} that fails DNS/TLS for all four. Read as a
single fact about the parent domain, this would suggest all four are
simply unreachable. Two of the four, CIFT and CIBA, are reachable under their
own \texttt{.res.in} domain instead (\texttt{cift.res.in},
\texttt{ciba.res.in}, both 200); the other two are not confirmed reachable
under any working domain (CIFRI has no \texttt{.res.in} equivalent probed
successfully here, and CIFA's \texttt{.nic.in} domain returns 400). The
correction is specific and only two-for-four: an institution sharing a dead
parent-domain pattern is worth a second, institution-specific probe before
being written off with the other three, but that second probe does not
always succeed.

\para{What changed since \ProbeDate{}} Re-probing \NEndpointsReprobed{} of
the \NSourcesSurveyed{} endpoints while writing this paper, \AsOf{}, finds
\NEndpointsUnchanged{} identical to the \ProbeDate{} record --- WCPFC's
\texttt{/public-domain} still 404, PANGAEA's OAI-PMH endpoint still 200 and
its old Elasticsearch endpoint still 404, \texttt{cift.res.in} and
\texttt{ciba.res.in} still 200, VertNet's appspot API still 500, and the
OBIS statistics call filtered to the IndOBIS node still returning exactly
\IndObisDatasets{} datasets, \IndObisRecords{} records, and
\IndObisSpecies{} species over \IndObisYearStart{}--\IndObisYearEnd{},
byte-identical to \ProbeDate{}'s figures. The remaining \NEndpointsDrifted{}
have not: \texttt{fishbase.ropensci.org}, dead by DNS resolution failure on
\ProbeDate{}, now resolves to a live host whose TLS handshake completes and
serves a self-signed certificate with subject and issuer both
``\FishbaseCertCN{}'' (\texttt{openssl s\_client}'s verify return code
\FishbaseOpenSSLVerifyCode{}, self-signed certificate) --- an ingress with no
route configured for this hostname, a different failure mode from an
unresolvable name, though the practical verdict (do not code against it) is
unchanged. VertNet's portal, 404 on \ProbeDate{}, now answers 308 and
redirects to a live 200 marketing page at \texttt{www.vertnet.org}, while its
programmatic \texttt{appspot.com} search API is still 500 --- the verdict
for programmatic use (defunct) is unchanged even though the specific status
code this survey's own table cites for ``portal'' no longer holds. Both
drifted cases were probed once more, from this session's own network access,
which is a real limitation on how confidently a single re-probe can be read
--- discussed in \S\ref{sec:limitations} --- but the FishBase case in
particular shows a protocol-level, server-side signal (a live TLS listener
answering with a specific, named default certificate) that a routing or
caching difference between vantage points does not plausibly explain.

\section{Discussion}
\label{sec:discussion}

The sharpest single finding in this survey is not any one dead or gated
source; it is what IndOBIS's status does to the otherwise clean statement
that no Indian marine data source has a machine-readable API. Read at the
level of institutions, that statement holds without exception: CMFRI, FSI,
CIFT, CIBA, CIFRI, and CIFA between them offer document repositories, portal
pages, and manual-request forms, and not one of the six runs a queryable API
of its own. IndOBIS looks, at first glance, like the seventh Indian
institution and the sole counterexample --- a Regional OBIS Node run by
CMLRE (the Centre for Marine Living Resources and Ecology), returning 200 on
its own domain (\texttt{indobis.in}) with a real dataset behind it. It is
not actually a counterexample, and the reason why is the finding worth
acting on: CMLRE has not built and does not run an API. IndOBIS is a
\emph{node designation inside the OBIS aggregator}, reachable by passing
\texttt{?nodeid=\IndObisNodeId{}} to the exact same
\texttt{/v3/occurrence}, \texttt{/v3/dataset}, and \texttt{/v3/statistics}
endpoints Table~\ref{tab:open} already lists as open and already implemented
in this codebase's \texttt{obis.py}. Filtering on that node returns
\IndObisDatasets{} datasets, \IndObisRecords{} records, and
\IndObisSpecies{} species across \IndObisYearStart{}--\IndObisYearEnd{}
(the Actinopterygii subset alone, \texttt{taxonid=10194}, is
\IndObisActinoRecords{} records across \IndObisSpecies{} $\to$
\IndObisActinoSpecies{} species and \IndObisActinoDatasets{} datasets), all
of which a call to \texttt{obis.py} already pulls in as part of the global
OBIS query it already makes.

The practical consequence is specific and easy to get wrong in the direction
of doing unnecessary work: the correct response to ``does this platform have
Indian marine occurrence data'' is not a new integration project aimed at
CMLRE, because there is no CMLRE-run endpoint to integrate with --- it is
adding the node filter to a query already being issued, as a
provenance/diagnostic facet (``how much of this region's background is
Indian-sourced'') rather than as a source of additional data. A survey that
stopped at ``IndOBIS returns 200, mark it reachable'' would have filed it
in the same bucket as CMFRI's eprints repository or CIFT's own domain ---
reachable, but requiring integration work proportional to a new source. The
finding that makes this survey worth having is the one level down from
reachability: \emph{why} it is reachable, which changes the item from a
backlog entry into a one-line filter on code that already ships.

\section{Limitations}
\label{sec:limitations}

This survey is the work of one team, probing largely on one date
(\ProbeDate{}), against each source's documented or discoverable endpoint
from this team's own network vantage point --- a different vantage point
(a different country, a different institutional network, a service behind
geofencing or a CDN edge) could in principle return a different status for
the same source at the same moment, and neither the original pass nor this
paper's re-verification pass records which network it was run from. The
three-way verdict (open, gated, dead) is a deliberate simplification that
several sources already strain: AIMS, DFO Canada, and NOAA Stock SMART all
show a documentation page or portal succeeding while the actual data
endpoint fails, which is neither cleanly ``open'' nor cleanly ``dead,'' and
is reported in Table~\ref{tab:gated}'s detail column rather than
folded into a fourth category this survey does not otherwise use. The
re-verification pass in \S\ref{sec:results} covers \NEndpointsReprobed{} of
\NSourcesSurveyed{} sources, not all of them, chosen for being the most
specific or surprising claims in the original pass rather than a random or
complete sample, so \NEndpointsUnchanged{} of \NEndpointsReprobed{} holding
steady is evidence about those particular endpoints and should not be read
as an estimate of how much of the full table has drifted.

None of this is presented as a flaw to correct in a future version, because
a future, ``final'' version of this survey is not a coherent goal. The
paper's own thesis is that this kind of survey decays on a timescale of
weeks, not years, which \S\ref{sec:results}'s own re-verification already
demonstrates for \NEndpointsDrifted{} of \NEndpointsReprobed{} endpoints
checked barely five weeks after the original pass; the right response is not
a more careful one-time survey but a habit of re-probing before relying on a
``works'' verdict months later, exactly as \texttt{docs/data-sources.md}'s
own header already tells its readers to do.

\section{Conclusion}

Of \NSourcesSurveyed{} marine biodiversity, fisheries, and catch/effort data
sources probed live on \ProbeDate{}, \NTierOne{} are open and
occurrence-shaped with no key required, \NTierTwo{} are reachable but gated
or bulk-download-only, and \NTierThree{} Indian institutions have, with one
exception, no machine-readable API at all. That exception is the survey's
sharpest finding: IndOBIS needs no new integration because it is already
reachable as a node filter on the OBIS API this codebase already calls, which
turns what looks like an integration backlog item into a one-line addition
to an existing query. Three further corrections --- that GBIF is not a
superset of OBIS for marine fish, that a single status-000 probe conflates a
transient flake with a genuinely dead host, and that a dead parent domain
does not mean a dead institution --- show that even a survey built from live
probes needs a second check before a verdict is trusted, the same discipline
this research programme has applied elsewhere to its own computed metrics.
Every claim in this paper is graded Dated rather than Artifact, because the
object under test is infrastructure outside this project's control, and a
partial re-verification performed while writing this paper already finds
two of nine re-checked endpoints answering with a different technical
signature than they did five weeks earlier --- direct evidence for, not
against, the paper's central argument that a survey like this is worth
publishing now, with a date attached, rather than waiting for a version that
this genre of survey cannot produce.
